\section{重尾分布：当均值与方差失效时}
\label{sec:05-Probability/06-Common-Distributions/07}

你从交易所抓了一批日收益率数据，前三十天一切正常：均值 \(0.05\%\)，标准差 \(1.2\%\)。两个数字算出来，顺手写进报告。第三十一天，市场跌了 \(8\%\)。你重新拉一下均值——掉到 \(-0.2\%\)。第三十二天又跌 \(6\%\)，均值继续往下跳。你不慌：大数定律说样本量越大均值越稳，等着就是了。

可如果你碰巧活在 Cauchy 的世界里，等多久都没用。

任何一批有限数据，你搜刮出来的样本均值和样本方差永远是有限数——给一列实数做加减乘除，不可能算出无穷。于是人很容易默认：背后的总体矩也一定存在，只是暂时还不知道具体数值。这个默认是错的。矩是积分，积分可以发散；上一节的 Student \(t\) 在自由度 \(\nu\le1\) 时均值已经不存在了。本节把这条线索展开到底：当极端值来得太频繁，频繁到平均无法将其稀释时，``均值''和``方差''这两个最常用的概括工具，会次第失效。

\subsection{Cauchy：一切正常的表象，发散的积分}

自由度 \(\nu=1\) 的 Student \(t\) 就是标准 Cauchy 分布，密度

\[
f(x)=\frac{1}{\pi(1+x^2)},\qquad x\in\R.
\]

对称，钟形，处处光滑——肉眼扫过去和 Gaussian 一样体面。

但当你试图计算均值：

\[
\E|X| = \int_{-\infty}^{\infty}\frac{|x|}{\pi(1+x^2)}\,dx = \infty.
\]

被积函数在远处的行为是 \(|x|/(\pi x^2)\sim 1/(\pi|x|)\)——衰减得不够快，面积发散到无穷。注意是绝对值期望发散，不是正负相消的问题：正半轴发散到 \(+\infty\)，负半轴发散到 \(-\infty\)，它们的差是一个 \(\infty-\infty\) 型的未定式，不能赋予任何有限值。均值不存在。

后果不是纸面上的数学奇观。设 \(X_1,X_2,\ldots\) 独立同分布于标准 Cauchy。一个令初学者不安的事实：

\[
\bar X_n = \frac1n\sum_{i=1}^n X_i \sim \text{Cauchy},
\]

不是近似，是精确成立。\(n=10\) 的样本均值和 \(n=10^6\) 的样本均值的波动幅度完全一样。你不断采样、不断取平均，均值永远稳定不下来。

为什么会这样？直觉上一个单独的新观测，有不可忽略的概率大到超过此前所有观测的累计量级——一项就能把均值重新拽走。这个直觉有精确的数学对应。矩的有限性和尾部概率的衰减速度由一条充要条件直接挂钩：

\[
\E|X|<\infty \iff \sum_{n=1}^{\infty}P(|X|>n)<\infty.
\]

Cauchy 让右端级数发散，尾部厚到概率加起来都爆掉。再用 Borel--Cantelli 引理一走：事件 \(\{|X_n|>n\}\) 会无限次发生。无论已经收集了多少数据，未来总会再来一个足够大的观测，大到单枪匹马推翻此前积累的全部平均。

\subsection{轻尾和重尾的分界线到底画在哪}

Cauchy 是理解重尾的窗口，不是用来猎奇的病态孤例。判断尾部轻重有两条递进标准：

\begin{itemize}
\tightlist
\item \textbf{矩标准}：\(\E[|X|^k]\) 对哪些 \(k\) 有限。Gaussian 和 Exponential 所有阶矩都有限；Cauchy 连一阶都过不去。
\item \textbf{指数矩标准}：\(\E e^{t|X|}\) 是否对某个 \(t>0\) 有限。这等价于问尾部至少按指数速度衰减，比``所有幂次矩有限''更强。
\end{itemize}

两条都通过的是轻尾。Gaussian 的尾部 \(P(X>x)\) 按 \(e^{-x^2/2}/x\) 衰减，Exponential 按 \(e^{-\lambda x}\) 衰减——极端值的概率被指数急速压低，平均不需要太多数据就把它们稀释干净。通不过指数矩标准的是重尾。

重尾最典型的形状是幂律。Pareto 分布的生存函数干脆就是幂函数：

\[
P(X>x)=\Bigl(\frac{x_m}{x}\Bigr)^{\alpha},\qquad x\ge x_m>0,
\]

参数 \(\alpha>0\) 叫尾指数。尾部按幂次衰减，比任何指数都慢——因为对任何 \(\lambda>0\)，\(x^{-\alpha}\) 最终一定被 \(e^{-\lambda x}\) 追上，但幂律耗的是之前漫长的大 \(x\) 区间。直接积分可知矩的存在条件极简洁：

\[
\E[X^k]<\infty \iff k<\alpha.
\]

尾指数像一架梯子，每爬过一个整数，就解锁一批概括工具：

{\def\LTcaptype{none}
\begin{longtable}{@{}ll@{}}
\toprule\noalign{}
尾指数范围 & 新变得可用的概括 \\
\midrule\noalign{}
\endhead
\bottomrule\noalign{}
\endlastfoot
\(\alpha\le1\) & 均值不存在，位置只能用中位数和分位数描述 \\
\(1<\alpha\le2\) & 均值 \(\E[X]\) 存在，大数定律可用 \\
\(2<\alpha\le3\) & 方差 \(\Var(X)\)，经典中心极限定理 \\
\(3<\alpha\le4\) & 三阶矩，偏度可定义 \\
\(\alpha>4\) & 四阶矩，峰度可定义 \\
\end{longtable}
}

Cauchy 正坐在 \(\alpha=1\) 的边界：\(P(|X|>x)\sim 2/(\pi x)\)，尾指数恰好为 1，属于均值不存在的那一档。

这个梯子还有一条``夹层''——对数正态分布。若 \(\log X\sim\mathcal N(\mu,\sigma^2)\)，所有幂次矩：

\[
\E[X^k] = \exp\Bigl(k\mu+\frac{k^2\sigma^2}{2}\Bigr) < \infty,
\]

但 \(\E e^{tX}=\infty\) 对一切 \(t>0\) 成立。按矩标准它温和，按指数矩标准它是重尾。两条标准画的不是同一条线。讨论``重尾''的时候，先声明用哪条标准，否则鸡同鸭讲。

\subsection{现实数据里的重尾}

以下领域的数据反复被验证呈幂律型重尾：

\begin{itemize}
\tightlist
\item 保险理赔与再保险损失：少数巨灾贡献大部分赔付额；
\item 金融资产收益：日收益绝对值的尾部常用 \(\alpha\approx 3\) 描述——方差勉强存在，四阶矩已经可疑；
\item 网络流量与文件大小：突发流量和自相似性的来源；
\item 城市人口、公司规模、词频：Zipf 型幂律，\(\alpha\) 接近 1。
\end{itemize}

给这类数据套 Gaussian 模型，错的不是中心，是尾巴——而风险正好住在尾巴里。

算一笔数量级对比。把两个分布都标准化——减去自己的均值、除以自己的标准差（如果存在）——然后问：``超出均值 5 个标准差''的概率是多少。Gaussian 回答：

\[
P(X>\mu+5\sigma) \approx 3\times 10^{-7},
\]

几百万分之一。尾指数 \(\alpha=3\) 的 Pareto 在同样的标准化下给出大约 \(5\times 10^{-3}\)，千分之五。差距上万倍。如果你的风险模型用 Gaussian 定价，几百万年一遇的损失会每年遇上好几次。这不仅是参数没估准——这是分布族选错了，问题出在结构层面。

\subsection{中心极限定理在重尾世界的残存范围}

方差有限时——对 Pareto 即 \(\alpha>2\)——经典中心极限定理的条件仍然满足：\(\sqrt{n}\) 缩放，Gaussian 极限成立。但 Berry--Esseen 误差界由三阶绝对矩控制（见\hyperref[sec:05-Probability/08-Limit-Theorems/03]{``中心极限定理与 Gaussian 近似''}）：重尾分布的三阶矩很大甚至不存在，收敛可以极慢。更要命的是，CLT 控制的是分布函数的绝对误差——中心区域已经近似得很好的时候，尾部相对误差仍可以是数量级的。用 Gaussian 近似回答``中心附近怎样''可以，回答``极端事件多频繁''不行。

方差无穷时——Pareto 的 \(\alpha\le2\)——经典 CLT 的条件整体失败。\(\sqrt{n}\) 缩放不再适用。对 \(\alpha<2\) 的规则变化尾部，在适当中心化和尾部平衡条件下，缩放的量级是 \(n^{1/\alpha}\)，极限对象不再是 Gaussian，而是 \(\alpha\)-稳定分布。Cauchy 正是 \(\alpha=1\) 的稳定分布——样本均值精确保持 Cauchy，就是稳定分布自相似性质的一个特例。边界 \(\alpha=2\) 需要单独处理：Pareto 型尾部常出现 \(\sqrt{n\log n}\) 量级的缓慢修正，不能直接把 \(\alpha=2\) 扔进 \(n^{1/\alpha}=\sqrt{n}\) 公式。

稳定分布的完整理论超出本单元。这里只需刻下一条认知：均值与方差的语言失效，不等于没有极限规律——只是极限规律可能换了缩放方式、换了极限对象。

\subsection{实务：先看尾，后选工具}

面对新数据，顺序应该反过来——别上来就求均值方差，先看一眼尾部。

矩不存在在数据里有可观察的表现。轻尾分布下，随着样本量增加，样本均值和样本方差会逐渐稳定；重尾下（\(\alpha<2\)），样本方差会被不断抵达的极端值一次一次地抬高，跑动图里它拒绝收敛——估计量本身拒绝稳定，就是对总体矩可能不存在的活体诊断。

具体操作顺序：

\begin{enumerate}
\tightlist
\item 画直方图；正数据可先取对数再画。看尾部是否远比 Gaussian 的薄壳厚。
\item 画经验尾部图：在对数--对数坐标下画 \(P(X>x)\) 的经验版本（样本中超过 \(x\) 的比例）。幂律会呈现近似直线，斜率约为 \(-\alpha\)。
\item 若尾部明显偏重，位置和离散度改用中位数、分位数、四分位距——它们对任意分布都存在，不被个别极端值绑架。
\item 确实需要量化尾部时再估计 \(\alpha\)。Hill 估计的思路：取样本中最大的若干个值，用它们相对于某个高门槛的平均对数超出量估计尾部衰减速度。门槛选择对结果极敏感，收敛慢，几百个样本给出的 \(\widehat\alpha\) 慎读——它常常只是一个有数字外表的猜测。
\end{enumerate}

本节的迁移要点不是``记住 Cauchy 的密度公式''，而是一个层级递进的判断结构：有限样本里算出来一个数，不能反推总体矩存在；矩如果不存在，以矩为基础的一切方法——均值、方差、OLS、CLT 近似——都踩在空气上。重尾世界里，诚实的概括从分位数开始，把均值留给那些配得上它的分布。
